\section{The physical Fourier cutoff does not intertwine the global PW projection}
\label{sec:physical-projection-obstruction}

\paragraph{Status and exact question.}
The complete author and root independent whole-proof readings
are recorded. Further independent review has its own event scope.
It tests a literal operator identity left open in
\texttt{MA-Q1-CQ003} and \texttt{MA-Q1-RF003}.  It does not
identify the local and global Hilbert norms, or replace a
projection by its trace.  The counterexample is first constructed
on genuine Schwartz PW data and then transferred to the original
coordinate-flat LF core of MA-GA-001.

Let $\chi=\chi_4$ be the primitive odd character modulo $4$,
and put $S=\{\infty,2\}$.  Write
\[
 \chi(n)=0\ (2\mid n),\qquad
 \chi(4n+1)=1,\qquad\chi(4n+3)=-1,
 \qquad \beta(z)=L(z,\chi).
 \tag{MA-PI.1}
\]
For an odd archimedean Schwartz function $f$ define the actual
centered global theta image in logarithmic coordinates by
\[
 Y_f(x)=2e^{x/2}\sum_{n\geq1}\chi(n)f(ne^x).
 \tag{MA-PI.2}
\]
This is the primitive character model of MA-CGP-001, with
the actual finite vector
$\theta_2=\chi|_{\mathbb Z_2^\times}\mathbf1_{\mathbb Z_2^\times}$
at $2$ and $\mathbf1_{\mathbb Z_p}$ outside $S$.
Both adelic polar moments vanish at the ramified finite place.
The image $Y_f$ belongs to the all-exponential Schwartz space
$\mathcal E$ for every such $f$.

Fix
\[
 0<a<1/2,\qquad s\geq0,\qquad
 H_{a,s}=L^2\!\left(\mathbb R,
           e^{2a|x|}(1+x^2)^{s/2}\,dx\right).
 \tag{MA-PI.3}
\]
Let $Q^{\mathrm{PW}}_{S,Y,a,s}$ be the actual orthogonal
projection onto the closed invariant global PW image of
product-norm Fourier level $Y$, with $S'=S$, in this same
Hilbert space.  This is the corrected centered realization
of the global projection in the specified MA-GP/MA-CGP/MA-GW
norm.  Masullo Definition 11.2.1.3, (11.2.1.7--8),
printed/PDF pp.~488--489, defines the corresponding
$a$-twisted level inside its $W^{s,2}$ realization; here
the image and orthogonal projection are explicitly fixed
in (MA-PI.3).  No inference that equivalent norms preserve
orthogonal projectors is made.  This is not the auxiliary
ordinary-frequency projector used to parametrize the range
in MA-GP/MA-CGP.

For each $Y>0$ let $\rho_Y\in C^\infty([0,\infty))$ satisfy
\[
 0\leq\rho_Y\leq1,\qquad
 \rho_Y(r)=1\quad(0\leq r\leq Y/2),\qquad
 \rho_Y(r)=0\quad(r\geq2Y).
 \tag{MA-PI.4}
\]
Every source-thin smooth PW-ball profile eventually has
these two plateau properties.  No derivative bound on the
family is required.  Use exactly the original physical
Fourier cutoff on $S$-local data:
\[
 \mathfrak B_{S,Y}
       =\mathcal F_S^{-1}M_{\rho_Y(|\xi|_S)}\mathcal F_S.
 \tag{MA-PI.5}
\]
Its actual descended version acts on the local quotient
completion constructed in MA-CQ and MA-GA.  The vectors
below lie in a canonical source subspace on which all
operations in (MA-PI.5) are literal Schwartz operations.

\paragraph{MA-PI-001 (an original flat-source obstruction at large finite levels).}
For each fixed $a,s$ in (MA-PI.3), there is $Y_0>0$ such
that, for every $Y\geq Y_0$ and every smooth profile
(MA-PI.4), there is an odd, archimedean-flat Schwartz
function $f_Y$ for which
\[
 Q^{\mathrm{PW}}_{S,Y,a,s}Y_{f_Y}
              \ne Y_{\mathfrak b_Y f_Y},
 \qquad
 \mathfrak b_Y f
     :=\mathcal F_\infty^{-1}
          \bigl(\rho_Y(4|\xi|)\widehat f(\xi)\bigr).
 \tag{MA-PI.6}
\]
The actual $S$-source input is
$\varphi_Y=f_Y\otimes\theta_2\in V_S$ of MA-GA.2.
Moreover
\[
 \mathfrak B_{S,Y}\varphi_Y
           =(\mathfrak b_Y f_Y)\otimes\theta_2
                         \in\mathcal S_0(\mathbb A_S).
 \tag{MA-PI.7}
\]
Thus (MA-PI.6) is a failure on the original flat core of
the canonical intertwining
\[
 Q^{\mathrm{PW}}_{S,Y,a,s}\,\vartheta_S
                 \stackrel{?}{=}\vartheta_S\,B_{S,Y},
 \tag{MA-PI.8}
\]
where $\vartheta_S$ is the actual centered global theta
embedding of the local coinvariant source and $B_{S,Y}$
is the descended physical cutoff.  Both sides are defined
on the witness.  The assertion holds at every narrow
width and every polynomial weight, including the draft's
$s\geq3$ scope, at arbitrarily large physical cutoff levels.
No zero-location assumption is used.

\begin{proof}
\emph{1. The actual finite source and the cutoff scale.}
Use the coherent additive characters of MA-NR-001/MA-CQ:
the archimedean Fourier kernel is $e^{-2\pi i x\xi}$ and
the $2$-adic kernel is its coherent positive fractional-part
counterpart.  The self-dual measure has
$\operatorname{vol}(\mathbb Z_2)=1$.
The two residue classes $1+4\mathbb Z_2$ and
$3+4\mathbb Z_2$ give directly
\[
 \mathcal F_2\theta_2(\xi)
     =\frac{i}{2}\chi(4\xi)
                           \mathbf1_{\mathbb Z_2^\times}(4\xi).
 \tag{MA-PI.9}
\]
Indeed on $\xi=u/4$, $u$ odd, their two integrals are
$\tfrac14(e^{\pi iu/2}-e^{3\pi iu/2})=\tfrac i2\chi(u)$;
at the other points of $\tfrac14\mathbb Z_2$ the difference
is zero, and outside that ball each integral is zero.
The finite Fourier support is exactly the shell
$|\xi|_2=4$.  Consequently multiplication by the physical
product-norm cutoff is exactly multiplication by
$\rho_Y(4|\xi_\infty|)$ on the archimedean Fourier factor.
Inverting both transforms proves (MA-PI.7), with the
finite phase canceled by the actual inverse.

The outside-$S$ integer balls restrict the global theta
sum to integers.  The two signs contribute equally because
both $\chi$ and $f$ are odd, giving (MA-PI.2) with no
extra constant.  The finite vector and its Fourier
transform both vanish at zero, so all these adelic
vectors are polar free independently of archimedean
moments.  The product formula makes the global theta
map constant on the $S$-unit coinvariant relations.
Thus the calculation is valid after taking the actual
local quotient; it is not a raw-cover representative
substituted for that quotient.

If $\widehat f$ is smooth, compactly supported, and has
support in $|\xi_\infty|\leq Y/4$, then the whole finite
Fourier support in (MA-PI.9) has product norm at most
$Y$.  It is compact and satisfies both polar conditions.
Its global theta image therefore belongs to the literal
invariant PW source level, and hence to
$\operatorname{Ran}Q^{\mathrm{PW}}_{S,Y,a,s}$.
Only this direct membership is used below; no equality
between different candidate PW ranges is needed.

\emph{2. An exact one-sided theta Gram form.}
For the intermediate calculation let $-1/2<b<1/2$ and
put $\sigma=1/2+b\in(0,1)$.
For $u\in C_c^\infty(\mathbb R_{>0})$, extend $u$ oddly
to $\mathbb R$ and put $f=\mathcal F_\infty^{-1}u$.
Define its logarithmic frequency vector by
\[
 A_u(r)=e^{(1/2-b)r}u(e^r),\qquad
 \widehat A_u(t)=\int_{\mathbb R}A_u(r)e^{-itr}\,dr.
 \tag{MA-PI.10}
\]
Each $A_u$ is a compact smooth function, and conversely
every compact smooth $A$ is obtained by
$u(\xi)=\xi^{b-1/2}A(\log\xi)$ on $\xi>0$.

Write $M_f(z)=\int_0^\infty f(t)t^{z-1}\,dt$.
The actual theta Mellin transform is
\[
 \mathcal M Y_f(z)=2\beta(z+1/2)M_f(z+1/2).
 \tag{MA-PI.11}
\]
It follows first by absolutely convergent unfolding in
$\Re(z+1/2)>1$ and then by analytic continuation; the
theta image is in $\mathcal E$ and $f$ is odd Schwartz.
The actual inverse odd Fourier transform gives
\[
 M_f(z)=i\,\frac{\Gamma_{\mathbb R}(z+1)}
                       {\Gamma_{\mathbb R}(2-z)}M_u(1-z),
 \qquad 0<\Re z<1.
 \tag{MA-PI.12}
\]
For clarity, $f(x)=2i\int_0^\infty u(r)\sin(2\pi xr)\,dr$.
Insert $e^{-\varepsilon x}$ in its Mellin integral and
integrate first in $x$.  The resulting gamma expression
converges uniformly for $r$ in the compact positive
support of $u$.  Letting $\varepsilon\downarrow0$ gives
$2i(2\pi)^{-z}\Gamma(z)\sin(\pi z/2)M_u(1-z)$.
The gamma reflection and duplication identities give
exactly (MA-PI.12).  This also fixes its Fourier sign.

Mellin Plancherel on the line $b$ now yields the exact
sesquilinear form
\[
\begin{split}
 \int_{\mathbb R}Y_f(x)\overline{Y_h(x)}e^{2bx}\,dx
   &=\frac1{2\pi}\int_{\mathbb R}
             W_b(t)\widehat A_u(t)
                         \overline{\widehat A_v(t)}\,dt,\\
 W_b(t)&=4\left|
       \frac{\Gamma_{\mathbb R}(\sigma+1+it)}
            {\Gamma_{\mathbb R}(2-\sigma-it)}
                         \right|^2|\beta(\sigma+it)|^2,
\end{split}
 \tag{MA-PI.13}
\]
where $h=\mathcal F_\infty^{-1}v$.
This is an inner product of actual theta images in a
specified one-sided norm.  It is not asserted to equal
the symmetric norm (MA-PI.3) at this stage.

The multiplier $W_b$ is a real, even, nonnegative tempered
function.  To verify the needed growth without a
mean-square theorem, let $A(X)=\sum_{n\leq X}\chi(n)$;
then $|A(X)|\leq1$.  Partial summation gives, for
$\Re z=\sigma>0$,
\[
 \beta(z)=\sum_{n\leq N}\chi(n)n^{-z}
               -A(N)N^{-z}+z\int_N^\infty A(x)x^{-z-1}\,dx,
 \tag{MA-PI.14}
\]
and the last two terms have absolute value at most
$N^{-\sigma}(1+|z|/\sigma)$.  Taking
$N\asymp1+|\Im z|$, and bounding the finite sum by
$\sum_{n\leq N}n^{-\sigma}$, proves
$\beta(\sigma+it)=O_\sigma((1+|t|)^{1-\sigma})$
for $0<\sigma<1$.  Stirling gives
\[
 \left|
       \frac{\Gamma_{\mathbb R}(\sigma+1+it)}
            {\Gamma_{\mathbb R}(2-\sigma-it)}\right|^2
       \sim (|t|/(2\pi))^{2\sigma-1}.
 \tag{MA-PI.15}
\]
Thus $W_b(t)=O_b(1+|t|)$ for each $|b|<1/2$.

\emph{3. Its inverse-Fourier distribution is not compactly supported.}
The precise functional equation needed here follows
from the same actual theta model.  Formula (MA-PI.9)
gives the finite theta phase $i$ and conductor $4$.
For $f_0(t)=t e^{-\pi t^2}$ one has
$\widehat f_0=-if_0$.  The actual Poisson identity
of MA-CGP.4 therefore reads
$Y_{f_0}(\log t)=Y_{f_0}(\log(1/(4t)))$.
Taking its Mellin transform and using
$2M_{f_0}(z)=\Gamma_{\mathbb R}(z+1)$ proves
\[
 \Gamma_{\mathbb R}(z+1)\beta(z)
       =4^{1/2-z}\Gamma_{\mathbb R}(2-z)\beta(1-z).
 \tag{MA-PI.16}
\]
Both sides continue as the same entire completed function.
Its entireness can also be seen directly from the rapid
theta decay at both ends and the entire reciprocal gamma
factor.  This proof records the root and conductor constants
instead of importing an unspecified functional-equation sign.

Put $d=1/2-b$.  Equations (MA-PI.13) and (MA-PI.16)
show that the real-axis multiplier has the entire continuation
\[
 \mathscr W_b(z)=4^{1-2b}\beta(d+iz)\beta(d-iz).
 \tag{MA-PI.17}
\]
Let $\kappa_b=\mathcal F_-^{-1}W_b$ in the tempered
distribution sense, with inverse normalization $1/(2\pi)$.
Suppose it had compact support in $[-R,R]$.  On a fixed
compact neighborhood of that support, a distribution
has finite order: for some $C,N$, its value is bounded
by $C$ times the maximum of derivatives through order
$N$ of the test function.  Applying this to a compact
cutoff times $e^{-izx}$ shows that its Fourier transform
is entire and satisfies
\[
 |\mathscr W_b(z)|\leq C'(1+|z|)^N e^{(R+1)|\Im z|}.
 \tag{MA-PI.18}
\]
Differentiation in $z$ is allowed in the same finite-order
bound, and agreement on the real axis identifies this
entire transform with (MA-PI.17).

This exponential-type bound is impossible.  Iterating
the gamma recurrence in (MA-PI.16) gives the exact values
\[
 \beta(-2n)=(-1)^n
        \frac{2\,4^n(2n)!}{\pi^{2n+1}}\beta(2n+1)
                  \qquad(n\geq0).
 \tag{MA-PI.19}
\]
For example the $n=0$ identity gives $\beta(0)=1/2$.
For the recurrence calculation one uses
$\Gamma(1/2-n)=(-4)^n n!\sqrt\pi/(2n)!$.
On the positive real axis the alternating series gives
$1-3^{-t}\leq\beta(t)\leq1$ for $t>0$.
At $T_n=2n+d$, (MA-PI.17) and (MA-PI.19) therefore give,
for $n\geq1$,
\[
 |\mathscr W_b(iT_n)|
    =4^{1-2b}|\beta(-2n)|\beta(2n+2d)
       \geq c_b(4/\pi^2)^n(2n)!,\qquad c_b>0.
 \tag{MA-PI.20}
\]
Since $(2n)!\geq n^n$, this grows faster than every
polynomial times $e^{C T_n}$.  It contradicts
(MA-PI.18).  Thus $\kappa_b$ is not compactly supported.
It is even because $W_b$ is even.

\emph{4. Separated genuine source data with nonzero Gram pairing.}
For any prescribed real cut positions $A<B$, there are
$A_u\in C_c^\infty((B,\infty))$ and
$A_v\in C_c^\infty(( -\infty,A))$ for which the
right side of (MA-PI.13) is nonzero.  We give the
distribution argument, including its quantifier.
If every such pairing vanished, then
$\kappa_b$ would vanish on $(-\infty,A-B)$.
Indeed take any compact test $\psi$ supported there,
choose $r_0>B$ with $r_0+\operatorname{supp}\psi\subset
(-\infty,A)$, approximate $\delta_{r_0}$ by smooth
compact functions in $(B,\infty)$, and pair their
convolutions with the fixed translated test
$\overline{\psi(\,\cdot-r_0)}$.
The assumed zero pairings pass to the distributional
limit and give $\langle\kappa_b,\psi\rangle=0$.
Evenness gives the corresponding vanishing on
$(B-A,\infty)$, which would make $\kappa_b$ compactly
supported.  Step 3 excludes this.

This argument supplies ordinary smooth compact functions,
not distributional input vectors: choose a nonzero test
of $\kappa_b$ beyond the prescribed gap, then take a
sufficiently concentrated member of the indicated
approximate identity.  The limiting pairing is nonzero,
so an actual finite member already has nonzero pairing.
Equivalently one may search any fixed countable dense
families of compact smooth functions in the two open
half-lines; continuity of (MA-PI.13) guarantees a pair
with nonzero absolutely convergent integral.  No claim
is made for every arbitrarily preassigned pair of bounded
bands.  The assertion is existence of suitable compact
bands beyond each prescribed physical cut and gap.

Use this result with $b=-a$, $A=\log(1/8)$ and
$B=\log(1/2)$.  Choose such $u,v$, and let
$f=\mathcal F_\infty^{-1}u$,
$h=\mathcal F_\infty^{-1}v$.  They are odd Schwartz,
with smooth compact Fourier supports respectively in
$|\xi|>1/2$ and $0<|\xi|<1/8$, and
\[
 D:=\int_{\mathbb R}Y_f(x)\overline{Y_h(x)}e^{-2ax}\,dx
                           \ne0.
 \tag{MA-PI.21}
\]
All functions here are actual primitive theta source
vectors and their actual images.  There is no requirement
that $f$ or $h$ be flat at zero at this intermediate step.

\emph{5. Large physical levels in the specified symmetric weighted norm.}
For $L>0$ use the unitary archimedean frequency dilation
\[
 u_L(\xi)=e^{-L/2}u(e^{-L}\xi),\qquad
 v_L(\xi)=e^{-L/2}v(e^{-L}\xi),\qquad Y=e^L.
 \tag{MA-PI.22}
\]
Their inverse transforms are
$f_L(t)=e^{L/2}f(e^Lt)$ and
$h_L(t)=e^{L/2}h(e^Lt)$.
Directly from the actual series (MA-PI.2),
\[
       Y_{f_L}(x)=Y_f(x+L),\qquad
       Y_{h_L}(x)=Y_h(x+L).
 \tag{MA-PI.23}
\]
The supports chosen in step 4 and (MA-PI.4) imply
\[
 \mathfrak b_Yf_L=0,\qquad
 \mathfrak b_Yh_L=h_L,\qquad
 Y_{h_L}\in\operatorname{Ran}Q^{\mathrm{PW}}_{S,Y,a,s}.
 \tag{MA-PI.24}
\]
For the last statement the finite product Fourier norm
is at most $4(e^L/8)=Y/2<Y$, so step 1 applies.
These equalities use only the fixed and zero plateaus,
not the values or derivatives of the profile in transition.

The pairing in the specified global norm obeys
\[
 \begin{split}
 &e^{-2aL}L^{-s}
       \langle Y_{f_L},Y_{h_L}\rangle_{H_{a,s}}\\
 &\quad=\int_{\mathbb R}Y_f(x)\overline{Y_h(x)}
       e^{2a|x-L|-2aL}
       \frac{(1+(x-L)^2)^{s/2}}{L^s}\,dx
             \longrightarrow D\ne0.
 \end{split}
 \tag{MA-PI.25}
\]
For $s=0$ the factor $L^{-s}$ is one.  Pointwise, the
two weight factors tend to $e^{-2ax}$ and one.
For $L\geq1$, they are bounded by
$C_s e^{2a|x|}(2+|x|)^s$.  The two actual theta images
are in $\mathcal E$, so dominated convergence is valid.
Consequently the pairing in (MA-PI.25) is nonzero for
every sufficiently large $L$, independently of the
permitted profile.  Orthogonality and (MA-PI.24) give
\[
 \left\langle
 Q^{\mathrm{PW}}_{S,Y,a,s}Y_{f_L}
              -Y_{\mathfrak b_Yf_L},\,Y_{h_L}
                       \right\rangle_{H_{a,s}}
       =\langle Y_{f_L},Y_{h_L}\rangle_{H_{a,s}}\ne0.
 \tag{MA-PI.26}
\]
This is already an actual PW-source counterexample.
It is obtained for all sufficiently large $Y=e^L$, not
by asserting a failure of a limiting scalar trace.

\emph{6. Transfer to the original coordinate-flat LF core.}
Fix one such finite $Y$, its profile, and the corresponding
$f_L,h_L$.  Choose an even $\omega\in C_c^\infty(\mathbb R)$
equal to one near zero.  Put
\[
 f_{L,\varepsilon}(t)
        =(1-\omega(t/\varepsilon))f_L(t),\qquad
 r_\varepsilon(t)=\omega(t/\varepsilon)f_L(t).
 \tag{MA-PI.27}
\]
Then $f_{L,\varepsilon}$ is odd and flat in a neighborhood
of zero, and its tensor with $\theta_2$ is in the actual
source $V_S$ of MA-GA.2.  The approximation is not
asserted in the topology of flat functions at this step;
the precise global and cutoff-output estimates follow.

Because $f_L$ is odd, write
$r_\varepsilon(t)=\varepsilon g_\varepsilon(t/\varepsilon)$,
where
$g_\varepsilon(y)=\omega(y)f_L(\varepsilon y)/\varepsilon$.
For $0<\varepsilon\leq1$ these are a bounded family in
odd compactly supported Schwartz space.  Indeed the
zeroth derivative is bounded by the first derivative
of $f_L$ near zero; higher derivatives acquire factors
$\varepsilon^{j-1}$, and the fixed compact cutoff
handles every Schwartz seminorm.
The actual theta map is continuous from that Schwartz
space to $\mathcal E$, by the direct sum at infinity
and its Poisson-reflected sum at zero, as in MA-CGP-001.
Scaling (MA-PI.2) gives
\[
 Y_{r_\varepsilon}(x)
       =\varepsilon^{3/2}
          Y_{g_\varepsilon}(x-\log\varepsilon),\qquad
 \|Y_{r_\varepsilon}\|_{H_{a,s}}
       \leq C\varepsilon^{3/2-a}
                    (1+|\log\varepsilon|)^{s/2}
                         \longrightarrow0.
 \tag{MA-PI.28}
\]
The estimate follows from the same weight inequality
as (MA-PI.25).  Its exponent is positive throughout
the asserted narrow scope.

There is also convergence after the literal smooth
physical cutoff.  Differentiating the Fourier integral
of $r_\varepsilon$ gives
\[
 \sup_\xi|\partial_\xi^j\widehat r_\varepsilon(\xi)|
      \leq C_j\varepsilon^{j+2}
                         \qquad(j\geq0).
 \tag{MA-PI.29}
\]
For the fixed finite $Y$ and profile, multiplying by
$\rho_Y(4|\xi|)$ gives compactly supported smooth
functions tending to zero in every smooth seminorm.
The inverse Fourier transforms therefore tend to zero
in Schwartz space.  Theta continuity then yields
\[
 \|Y_{\mathfrak b_Yr_\varepsilon}\|_{H_{a,s}}
                   \longrightarrow0.
 \tag{MA-PI.30}
\]
No uniform derivative bound in the profile was used:
the flat approximation is selected after that profile
and finite level are fixed.

Since $\mathfrak b_Yf_L=0$, (MA-PI.28--30) and boundedness
of the actual orthogonal projection show that
\[
 \begin{split}
 &\left\langle
 Q^{\mathrm{PW}}_{S,Y,a,s}Y_{f_{L,\varepsilon}}
       -Y_{\mathfrak b_Yf_{L,\varepsilon}},\,Y_{h_L}
                  \right\rangle_{H_{a,s}}\\
 &\hspace{4em}\longrightarrow
                  \langle Y_{f_L},Y_{h_L}\rangle_{H_{a,s}}
                         \ne0.
 \end{split}
 \tag{MA-PI.31}
\]
Choose any sufficiently small positive $\varepsilon$.
This supplies the actual finite-$\varepsilon$ flat
witness $f_Y=f_{L,\varepsilon}$ in (MA-PI.6).
All original maps on it have the source and target
declared above; no distribution or nonflat PW vector
is substituted for the final source input.
\end{proof}

\paragraph{Exact source and conclusion scope.}
The invariant product-norm PW data and their polar
conditions are Masullo Definition 7.3.2.3,
(7.3.2.12--17), printed/PDF pp.~354--355, with $S'=S$.
The source global projection is Definition 11.2.1.3,
(11.2.1.7--8), pp.~488--489; the corrected centered
Hilbert realization used here is specified after
(MA-PI.3).  The physical local Fourier
cutoff is precisely the first bullet of Step II on
p.~516 and the composition on p.~517, in the proof of
Theorem 12.4.1.1.  These passages were freshly read for
the map tested here.  The coherent finite sign is the
explicit repair already fixed in MA-NR/MA-CQ; (MA-PI.9)
recomputes its required value directly.

The gamma form in (MA-PI.12) agrees with Tate (1950),
Theorem 2.4.1 and the real table on thesis p.~(2.14),
PDF p.~21.  The exact functional equation used to
prove the new obstruction is derived in (MA-PI.16),
and its required growth consequence is proved, not
imported from a mean-square or universality theorem.

The result refutes the canonical theta intertwining
(MA-PI.8), including on the actual flat core and with
source-thin smooth profiles at large finite levels.
It does not assert that the draft printed that exact
intertwining equation: its Step III kernel identification
does not itself supply such an equation or its proof.
The result explains why an actual quotient section and
equality of PW ranges cannot supply this literal map
identification by themselves.

The Gram-corrected global projection and the two actual
Gram factors in MA-GP-008/MA-CGP-003 remain valid.
The theorem neither rules out those different maps nor
changes the proved source-return operators.  It does
not refute an $A$-supported scalar trace identity,
remove the arithmetic FT002 divisor, or claim that
every conceivable factorization has failed.  Its
obstruction is a nonzero inner product of actual
source images in the specified global norm, witnessed
before any trace is taken.

\paragraph{Grothendieck placement.}
The character $\chi_4$ is the rank-one finite Artin
representation of $\operatorname{Gal}(\mathbb Q(i)/\mathbb Q)$
which is $-1$ on complex conjugation.  On
$\operatorname{Spec}\mathbb Z[1/2]$ it is the associated
finite-monodromy etale local system; its unramified Euler
polynomial is $1-\chi_4(p)T$, and its inertia-invariant
space at $2$ is zero.  The exact source conventions are
SGA~1, Expos\'e~V, \S7, and Deligne, \emph{Valeurs de
fonctions L}, Definition~6.1 and \S\S6.2--6.3, already
audited in MA-CH/MA-CQ and not claimed freshly reread here.
The norm and orthogonal projection in (MA-PI.3--8) are
analytic metric data; the obstruction does not alter
this Artin representation or assert a new weight or
functoriality statement.  The inherited cote~1
trace/determinant programme, BR-TR-01/02, scans~079--080
and~111--112, requires actual comparison maps before
trace cyclicity is applied.  Those archive locators are
inherited; the present proof determines one such proposed
map identity to be false, while retaining the actual Gram
and nuclear-return comparisons already constructed.
